% BHU PYQ -- Philosophy / Multidisciplinary: Adhyatma Vijnan (Sem III Mid Term 2025-26 NEP)
\documentclass[11pt,a4paper]{article}
\usepackage[a4paper,top=2.5cm,bottom=2.5cm,left=2.8cm,right=2.8cm,headheight=14pt]{geometry}
\usepackage{amsmath,amssymb}
\usepackage{fontspec}
\setmainfont{Kohinoor Devanagari}
\usepackage{microtype,fancyhdr,enumitem,hyperref,lastpage,parskip}

\hypersetup{colorlinks=true,urlcolor=black,linkcolor=black,pdftitle={PHIMD31: Adhyatma Vijnan -- BHU NEP 2025-26 Mid Term}}

\pagestyle{fancy}\fancyhf{}
\renewcommand{\headrulewidth}{0.4pt}\renewcommand{\footrulewidth}{0.4pt}
\fancyhead[L]{\small \textit{Banaras Hindu University}}
\fancyhead[R]{\small \textit{PHIMD31: Adhyatma Vijnan (Mid Term 2025-26)}}
\fancyfoot[C]{\small Page \thepage\ of \pageref{LastPage}}

\newlist{parts}{enumerate}{2}
\setlist[parts,1]{label=(\alph*),leftmargin=*,itemsep=4pt,topsep=3pt}
\setlist[parts,2]{label=(\roman*),leftmargin=*,itemsep=2pt,topsep=2pt}

\newcommand{\pts}[1]{\hfill \textit{[#1]}}

\begin{document}

\begin{center}
  {\large\bfseries BANARAS HINDU UNIVERSITY}\\[2pt]
  {\normalsize B. A. / B. Sc. (Hons.) Semester III Mid-Term Examination, 2025-26 (NEP)}\\[6pt]
  \rule{0.8\linewidth}{0.5pt}\\[6pt]
  {\large\bfseries MULTIDISCIPLINARY COURSE (PHILOSOPHY)}\\[4pt]
  {\large\bfseries Paper Code: PHIMD-31 / MD-31 : Adhyatma Vijnan (अध्यात्म विज्ञान)}\\[4pt]
  {\normalsize Time: 1:30 Hours \hfill Full Marks: 25}
\end{center}
\rule{\linewidth}{0.4pt}
\smallskip

\section*{SECTION -- I / दीर्घ उत्तरीय प्रश्न}
\noindent \textit{Note: Answer any \textbf{two} questions. Each question carries 10 marks.}\pts{2 $\times$ 10 = 20}

\noindent \textit{नीचे दिए गए प्रश्नों में से कोई दो प्रश्न कीजिए। प्रत्येक प्रश्न 10 अंकों का है।}

\bigskip

\begin{enumerate}[label=\textbf{\arabic*.},leftmargin=*,itemsep=14pt]

  \item अध्यात्म विज्ञान क्या है? अध्यात्म विज्ञान में दर्शनों के महत्त्व का विवेचन कीजिए।\pts{10}
  
  What is Adhyatma Vijnan? Discuss the importance of philosophy in Adhyatma Vijnan.

  \item पातञ्जल योग में प्रतिपादित अष्टाङ्ग योग का विवेचन कीजिए।\pts{10}
  
  Explain the Ashtanga Yoga as described in Patanjal Yoga.

  \item सांख्य दर्शन का संक्षिप्त परिचय देते हुए सत्कार्यवाद का वर्णन कीजिए।\pts{10}
  
  Describe the theory of Satkaryavada, while giving a brief introduction to Sankhya philosophy.

  \item गीता के पहले अध्याय में अर्जुन का दुख (विषाद) हमें क्या सिखाता है? उसका सन्देश आज के जीवन में कैसे उपयोगी है, इसका वर्णन करें।\pts{10}
  
  What does Arjuna's sorrow (Vishada) in the first chapter of the Gita teach us? How is its message useful in today's life?

\end{enumerate}

\bigskip
\rule{\linewidth}{0.4pt}
\bigskip

\section*{SECTION -- II / लघु उत्तरीय प्रश्न}
\noindent \textit{Note: Answer all \textbf{five} short questions. Each question carries 1 mark.}\pts{5 $\times$ 1 = 5}

\noindent \textit{नीचे दिए गए लघु प्रश्नों का समाधान कीजिए। प्रत्येक प्रश्न 1 अंक का है।}

\bigskip

\begin{enumerate}[label=\textbf{\arabic*.},leftmargin=*,itemsep=10pt]

  \item गीता के त्रयोदश (13) अध्याय का नाम क्या है?\pts{1}
  
  What is the name of the thirteenth chapter of Bhagavad Gita?

  \item भगवान् श्रीकृष्ण के शङ्ख का नाम क्या है?\pts{1}
  
  Mention the name of the conch (Sankha) of Lord Krishna.

  \item बृहदारण्यकोपनिषत् किस वेद से सम्बद्ध है?\pts{1}
  
  Brihadaranyakopanishad belongs to which Veda?

  \item उत्तर मीमांसा दर्शन कितने प्रकार के प्रमाणों को स्वीकार करते हैं?\pts{1}
  
  How many types of proofs (Pramanas) are accepted by the Uttara Mimansa Philosophy?

  \item असत्कार्यवाद किस दर्शन का सिद्धान्त है?\pts{1}
  
  Which philosophy accepts Asatkaryavada?

\end{enumerate}

\end{document}
